\documentclass[a4paper]{article}

\usepackage[fontset=fandol]{ctex}
\usepackage{amsmath, amssymb}
\usepackage{graphicx}
\usepackage[margin=2.5cm]{geometry}
\usepackage[most]{tcolorbox}
\usepackage{hyperref}
\usepackage{booktabs}
\usepackage{float}
\usepackage{array}
\usepackage{xcolor}

\setlength{\parindent}{2em}
\setlength{\parskip}{0.35em}
\sloppy

\newcolumntype{L}[1]{>{\raggedright\arraybackslash}p{#1}}

\newtcolorbox{knowledgebox}[1]{
    enhanced, colback=blue!5!white, colframe=blue!75!black, colbacktitle=blue!75!black,
    coltitle=white, fonttitle=\bfseries, title={#1}, attach boxed title to top left={yshift=-2mm, xshift=2mm},
    boxrule=1pt, sharp corners
}

\newtcolorbox{importantbox}[1]{
    enhanced, colback=yellow!10!white, colframe=yellow!80!black, colbacktitle=yellow!80!black,
    coltitle=black, fonttitle=\bfseries, title={#1}, sharp corners
}

\newtcolorbox{warningbox}[1]{
    enhanced, colback=red!5!white, colframe=red!75!black, colbacktitle=red!75!black,
    coltitle=white, fonttitle=\bfseries, title={#1}, sharp corners
}

\newcommand{\notetitle}{各式各样的 Self-attention 变型}
\newcommand{\notesubtitle}{从稀疏化、低秩压缩到 Linear Attention 与 Synthesizer}
\newcommand{\noteauthors}{基于李宏毅老师《机器学习》课程整理}
\newcommand{\notedate}{\today}
\newcommand{\videochannel}{李宏毅机器学习深度学习课程（Bilibili）}
\newcommand{\videoduration}{约 71 分钟}
\newcommand{\videourl}{https://www.bilibili.com/video/BV1TAtwzTE1S?p=17}

\begin{document}
\pagestyle{empty}

\begin{center}
    \vspace*{1.4cm}
    {\Huge\bfseries \notetitle\par}
    \vspace{0.35cm}
    {\LARGE \notesubtitle\par}
    \vspace{0.8cm}
    \includegraphics[width=0.62\textwidth]{video.jpg}\par
    \vspace{0.8cm}
    {\large \noteauthors\par}
    {\large \notedate\par}
    \vspace{0.8cm}
    \begin{tcolorbox}[width=0.84\textwidth,center,sharp corners,
                      colback=black!3!white,colframe=black!50!white]
        \centering
        \textbf{视频信息}\\[3pt]
        频道：\videochannel \\
        时长：\videoduration \\
        链接：\href{\videourl}{\videourl}
    \end{tcolorbox}
\end{center}

\newpage
\tableofcontents
\newpage
\pagestyle{plain}

% =====================================================================
\section{为什么要改造 Self-attention}
% =====================================================================

标准 self-attention 的优势是每一个位置都可以直接读取整个序列的信息。若输入序列为
\[
    A=(a^1,a^2,\ldots,a^N),
\]
模型会为每个位置产生 query、key 与 value：
\[
    q^i=W^q a^i,\qquad k^i=W^k a^i,\qquad v^i=W^v a^i.
\]
接着每个 query 都要和每个 key 做相似度计算。若不考虑多头和缩放常数，一个常见写法是
\[
    e_{ij}=q_i^\top k_j,\qquad
    \alpha_{ij}=\frac{\exp(e_{ij})}{\sum_{\ell=1}^{N}\exp(e_{i\ell})},
    \qquad
    b^i=\sum_{j=1}^{N}\alpha_{ij}v^j.
\]
这里真正让人头痛的是 $\alpha\in\mathbb{R}^{N\times N}$。序列长度是 $N$，attention matrix 就有 $N^2$ 个位置。对短句子来说这还可以接受；对长文本、语音、图像 patch、蛋白质序列、长上下文语言模型来说，$N^2$ 很快会成为主要瓶颈。

\begin{figure}[H]
    \centering
    \includegraphics[width=0.86\textwidth]{figures/fig01_attention_matrix_cost.jpg}
    \caption{标准 self-attention 需要产生 $N\times N$ 的 attention matrix；每个 query 都要和每个 key 计算一次相似度，因此长度一大，矩阵本身就成为计算和记忆体瓶颈。\protect\footnotemark}
    \label{fig:attention-cost}
\end{figure}
\footnotetext{视频画面时间区间：约 02:15--03:05；字幕附近说明「sequence 长度是 $N$」「有 $N$ 个 key、有 $N$ 个 query」「得到一个 $N$ 乘以 $N$ 的矩阵」。}

复杂度可以粗略分成两块。第一块是产生 $Q,K,V$ 的线性变换，复杂度大约随 $N$ 线性成长；第二块是 $QK^\top$ 与 attention matrix 乘 value 的过程，会出现 $N^2$。当 $N$ 不大时，前者和后者未必差很多；当 $N$ 极大时，后者通常会主导。

\begin{knowledgebox}{本讲的共同问题}
    这堂课讨论的各种方法看似五花八门，其实大多在回答同一个问题：我们能不能不要真的算完整的 $N\times N$ attention matrix？有的方法让矩阵变稀疏，有的方法让矩阵变瘦，有的方法完全改变乘法顺序，还有的方法甚至不再由 $Q$ 和 $K$ 产生 attention 权重。
\end{knowledgebox}

\subsection{改造 self-attention 的三条路线}

从课程内容看，可以把方法分成三类。

第一类是\textbf{稀疏化}：保留 $N\times N$ 这个概念，但只计算其中一小部分位置。其余位置被视为零，或被 mask 掉。Local attention、Longformer、BigBird、Reformer、Routing Transformer、Sinkhorn Sorting Network 都可以放在这个大方向下理解。

第二类是\textbf{压缩矩阵大小}：不再用所有 key 或 value，而是选出代表性的 key/value，或者用线性组合把长序列压成短序列。Linformer 和 compressed attention 代表这类思路。

第三类是\textbf{改变公式}：把 attention 看成矩阵乘法，尝试改变乘法顺序，或者用特征映射 $\phi(\cdot)$ 近似 softmax kernel，从而绕开显式 $N\times N$ 矩阵。Linear Transformer、Efficient Attention、Performer 这类方法都属于这一支。

\subsection{本章小结}

Self-attention 的强大来自全局两两交互，瓶颈也来自全局两两交互。若输入长度为 $N$，完整 attention matrix 有 $N^2$ 个元素；本讲所有变型都在尝试用结构假设、数据相关估计、低秩压缩或代数重写来减少这部分成本。

% =====================================================================
\section{用人类知识决定哪些位置不用算}
% =====================================================================

最直觉的加速方式是：有些位置明显不重要，干脆不要算。也就是说，我们先用人类对任务的理解设计一个 mask，只保留某些 query-key 对，其他位置直接设为零。这样做不需要模型先判断哪些位置重要，规则在训练前就固定。

\subsection{Local Attention 与 Truncated Attention}

Local attention 的核心假设是：很多任务中，一个位置最需要看的只是附近的位置。举例来说，语音或影像中相邻时间点、相邻 patch 的相关性往往较强；若只让每个位置看左右固定窗口，就能把 dense attention matrix 变成靠近对角线的带状矩阵。

\begin{figure}[H]
    \centering
    \includegraphics[width=0.82\textwidth]{figures/fig02_local_attention.jpg}
    \caption{Local attention 只计算对角线附近的蓝色区域，灰色位置直接设为零；这样可以减少 attention matrix 中需要计算的位置数。\protect\footnotemark}
    \label{fig:local-attention}
\end{figure}
\footnotetext{视频画面时间区间：约 08:17--09:11；字幕附近说明灰色位置「直接设为零」、只计算蓝色部分，并指出 local/truncated attention 只能看小范围资讯。}

若窗口大小为 $w$，每个 query 只看大约 $w$ 个 key，复杂度可由 $O(N^2)$ 降为 $O(Nw)$。当 $w$ 远小于 $N$ 时，这个降幅很可观。

但 local attention 的缺点也很清楚：它牺牲了 self-attention 最吸引人的性质，也就是任意两个位置可以一步交互。若某个位置需要远距离信息，local attention 必须靠堆叠很多层，让信息一层一层传过来。这样它会变得有点像 CNN：局部感受野便宜，但长距离依赖要靠层数累积。

\begin{warningbox}{Local attention 的代价}
    Local attention 不是免费午餐。它减少计算，是因为它明确禁止某些远距离连接。若任务强烈依赖远距离关系，例如长文本的跨段指代、长音频中的全局语义、图像中相隔很远区域的对应关系，固定局部窗口可能会漏掉重要信息。
\end{warningbox}

\subsection{Global Attention：给少数 token 全局权限}

为了弥补 local attention 的局部性，可以加入少数 global token。做法是指定一些特殊位置，让它们可以 attend 到所有 token，也允许所有 token attend 到它们。这样全局信息可以先汇聚到 special token，再由 special token 分发给其他位置。

\begin{figure}[H]
    \centering
    \includegraphics[width=0.86\textwidth]{figures/fig03_global_attention_token.jpg}
    \caption{Global attention 会让特殊 token 读取所有位置，也让所有位置读取特殊 token；这些 token 像全局信息交换站。\protect\footnotemark}
    \label{fig:global-attention}
\end{figure}
\footnotetext{视频画面时间区间：约 12:31--13:28；字幕附近说明 special token 会 attend 到所有 token，也会被所有 token attend，并提到可直接把序列开头 token 当成 special token。}

在文本任务中，BERT 的 CLS token 就很容易被想成一种全局汇聚点。对于分类任务，CLS 可以收集整句信息；对于长序列模型，少数 global token 可以在固定稀疏图上建立全局捷径。

不过，global token 也不是万能的。若所有长距离信息都必须经过少数几个节点，模型容量可能受到限制；如果 global token 选择不好，重要信息未必能有效通过它们传递。因此许多模型会把 local、global、random 或 dilated pattern 混合使用。

\subsection{Longformer 与 BigBird：固定稀疏图的组合}

Longformer 和 BigBird 展示了固定稀疏 pattern 可以有很多设计。Longformer 常见元素包括 sliding window attention、dilated sliding window，以及 global sliding window。BigBird 则在 local 和 global 之外加入 random attention，也就是随机允许一些远距离位置彼此连接。

\begin{figure}[H]
    \centering
    \includegraphics[width=0.92\textwidth]{figures/fig04_longformer_bigbird_patterns.jpg}
    \caption{Longformer 与 BigBird 使用不同的稀疏 attention pattern：局部窗口、扩张窗口、全局连接、随机连接可以组合起来，形成比纯 local 更灵活的稀疏图。\protect\footnotemark}
    \label{fig:longformer-bigbird}
\end{figure}
\footnotetext{视频画面时间区间：约 18:01--18:40；字幕附近说明 BigBird 包含 Longformer 的 local/global attention，并额外加入 random attention，论文中还有理论推导。}

随机连接听起来随意，但它有一个重要作用：在稀疏图中缩短任意两个位置之间的路径。纯 local attention 中，信息从序列左端传到右端要经过很多层；加入少量随机边后，图的直径可能显著下降。换句话说，random attention 不一定是为了精确建模某个语义关系，而是为了让信息传播更容易。

\begin{importantbox}{固定稀疏 pattern 的优缺点}
    固定 pattern 的优点是简单、可控、容易实现高效 kernel；缺点是它把「哪些位置重要」写死了。若任务中的关键依赖关系不符合人类设计的 pattern，模型可能会为了省计算而错过真正重要的连接。
\end{importantbox}

\subsection{本章小结}

用人类知识设计稀疏 pattern，是改造 self-attention 的第一条路。Local attention 用局部性换速度，global token 提供全局汇聚，Longformer 和 BigBird 则把多种连接方式组合起来。它们共同的特点是：哪些位置要算，在模型看见具体输入前就已经决定。

% =====================================================================
\section{让模型判断关键位置}
% =====================================================================

固定稀疏 pattern 很便宜，但未必最适合每个输入。于是第二个想法出现了：能不能让模型或算法根据当前输入判断哪些 attention 值可能大，只计算这些关键位置？

\subsection{只保留可能较大的 attention 值}

如果某些 attention weight 本来就接近零，那么把它们直接设为零，对最终 weighted sum 的影响可能很小。此时加速的关键不是「强行固定某种 pattern」，而是快速估计哪些位置可能重要。

\begin{figure}[H]
    \centering
    \includegraphics[width=0.82\textwidth]{figures/fig05_critical_parts_sparse.jpg}
    \caption{只关注 critical parts 的想法：深色块表示可能较大的 attention 值，浅色块可直接设为零，以减少实际计算。\protect\footnotemark}
    \label{fig:critical-parts}
\end{figure}
\footnotetext{视频画面时间区间：约 19:17--20:27；字幕附近说明若某些 attention weight 很小，可以直接设为零，并追问如何估计哪些位置可能有较大值。}

这个思路有一个微妙点：标准 attention 是先算出所有 $e_{ij}$，再通过 softmax 得到 $\alpha_{ij}$。如果你已经把所有 $e_{ij}$ 都算完，再丢掉小值，那计算量并没有省下多少。真正要加速，必须在完整 attention 计算之前，就用更便宜的方法筛出候选位置。

\subsection{Reformer 与 Routing Transformer：用 clustering 筛连接}

Reformer 和 Routing Transformer 采用相近思想：先把 query 和 key 按相似性分群，只让同一群内的 query-key 对计算 attention，跨群的连接则设为零。

\begin{figure}[H]
    \centering
    \includegraphics[width=0.86\textwidth]{figures/fig06_clustering_reformer.jpg}
    \caption{Clustering 方法会把相近的 query/key 分到同一群；只有同一 cluster 内的位置计算 attention，不同 cluster 的位置直接设为零。\protect\footnotemark}
    \label{fig:clustering}
\end{figure}
\footnotetext{视频画面时间区间：约 22:48--23:43；字幕附近说明红色 query 只与同 cluster 的 key 计算 attention，其余位置设为零，因此矩阵中很多位置不需要计算。}

这种方法背后的假设是：若 query 和 key 在表示空间中相距很远，它们的 dot product 往往不会大，最终 attention weight 可能很小。与其对所有位置精算，不如先用近似、较快的 clustering 找出可能相关的区域。

要注意的是，clustering 本身也有成本。若分群算法也需要 $O(N^2)$，那就失去了意义。因此实际方法会使用近似分群、哈希、routing 等技巧，让筛选阶段比完整 attention 更便宜。

\subsection{Sinkhorn Sorting Network：可学习的稀疏 pattern}

还有一种更大胆的方法：让另一个网络决定哪些格子要算。Sinkhorn Sorting Network 的图中，输入序列通过一个小网络产生类似 mask 的矩阵；这个 mask 决定 attention matrix 中哪些位置保留，哪些跳过。

\begin{figure}[H]
    \centering
    \includegraphics[width=0.88\textwidth]{figures/fig07_learnable_patterns_sinkhorn.jpg}
    \caption{Sinkhorn Sorting Network 用另一个可学习模块决定哪些 attention 位置要保留；图中的 mask pattern 与主网络一起训练。\protect\footnotemark}
    \label{fig:sinkhorn}
\end{figure}
\footnotetext{视频画面时间区间：约 29:46--30:46；字幕附近说明由网络产生较低解析度的 mask，例如先产生 $10\times10$，再放大成 $100\times100$，并可与整个网络一起反向传播训练。}

这里的困难是：是否计算某个格子本质上是 binary decision，但神经网络输出通常是 continuous。若直接把连续值硬切成 0 或 1，就可能造成不可微的问题。Sinkhorn 类方法的重点之一，就是用特殊技巧让「近似排序」或「近似二值 mask」仍可放进可训练系统。

\begin{knowledgebox}{从固定 pattern 到学习 pattern}
    Local attention 和 BigBird 这类方法把稀疏图写在架构里；Reformer 这类方法用输入相关的 clustering 产生稀疏图；Sinkhorn 进一步把稀疏图本身交给可学习模块决定。越往后，适应性越强，但实现和训练也越复杂。
\end{knowledgebox}

\subsection{本章小结}

数据相关稀疏化试图避免固定规则的僵硬性。它的共同目标是用较低成本预测哪些 attention 值值得精算。Reformer/Routing Transformer 通过 clustering 找相关区域；Sinkhorn Sorting Network 则让网络学习稀疏 pattern。关键风险在于：筛选过程本身不能太贵，也不能误删真正重要的连接。

% =====================================================================
\section{低秩与代表性 Key/Value}
% =====================================================================

前面的稀疏方法仍然保留 $N\times N$ attention matrix 的坐标系，只是不计算某些格子。另一条路线是问：我们真的需要完整的 $N$ 个 key 和 $N$ 个 value 吗？如果 attention matrix 里很多列是重复的、相依的，或可以由少数列线性组合得到，那么也许只需保留少数代表性 key/value。

\subsection{Attention Matrix 可能是低秩的}

Linformer 的直觉是：实际训练出来的 attention matrix 常常有冗余。若很多 column 之间高度相关，完整 $N\times N$ 矩阵包含的信息并没有看起来那么多。

\begin{figure}[H]
    \centering
    \includegraphics[width=0.84\textwidth]{figures/fig08_low_rank_attention.jpg}
    \caption{若 attention matrix 中许多 column 是冗余的或可由其他 column 线性组合得到，就可以把完整矩阵近似为较低秩的矩阵。\protect\footnotemark}
    \label{fig:low-rank}
\end{figure}
\footnotetext{视频画面时间区间：约 31:32--32:13；字幕附近说明 attention matrix 可能是 low rank，很多 column 可能是 duplicate 或 linear combination，因此不一定需要完整 $N\times N$。}

低秩观点可以写成矩阵近似。若原 attention matrix 为 $A\in\mathbb{R}^{N\times N}$，低秩假设大致认为它可以由较小的中间维度 $K$ 表示：
\[
    A \approx U V^\top,\qquad U,V\in\mathbb{R}^{N\times K},\qquad K\ll N.
\]
此时存储和计算都可以围绕 $K$ 展开，而不是围绕 $N$ 的平方展开。

\subsection{减少 Key 与 Value 的数量}

一个直接实现是挑出 $K$ 个代表性 key。原本每个 query 要和 $N$ 个 key 比较；若只和 $K$ 个代表性 key 比较，attention matrix 就从 $N\times N$ 变成 $N\times K$。对应地，value 也要压成 $K$ 个代表性 value，因为最后 weighted sum 需要和 key 一一对应的 value。

\begin{figure}[H]
    \centering
    \includegraphics[width=0.86\textwidth]{figures/fig09_representative_keys_values.jpg}
    \caption{只保留 $K$ 个代表性 key/value 后，attention matrix 由 $N\times N$ 变为 $N\times K$，每个 query 只需要和代表性 key 比较。\protect\footnotemark}
    \label{fig:representative-kv}
\end{figure}
\footnotetext{视频画面时间区间：约 32:49--33:46；字幕附近说明只挑有代表性的 key，矩阵会「变瘦」，并且 value 也要挑出对应的 $K$ 个代表性向量。}

这样做的优势是输出长度仍为 $N$：每个 query 都还在，所以每个输入位置仍能得到一个输出。它适用于需要对每个位置输出结果的任务，例如逐帧分类、序列标注、语音帧级判断等。

\subsection{能不能减少 Query？}

如果减少 query，情况就不一样了。query 的数量决定输出位置数：$N$ 个 query 会产生 $N$ 个输出；若只留下 $K$ 个 query，输出长度就变成 $K$。这在某些任务中可以接受，例如整段输入只需要输出一个分类标签；但在每个位置都要预测标签的任务中，减少 query 会直接改变输出形状。

\begin{figure}[H]
    \centering
    \includegraphics[width=0.86\textwidth]{figures/fig10_reduce_queries_caveat.jpg}
    \caption{减少 query 会改变输出序列长度；若任务需要每个输入位置都有一个 label，就不能随便把 query 数量变少。\protect\footnotemark}
    \label{fig:reduce-query}
\end{figure}
\footnotetext{视频画面时间区间：约 35:19--36:12；字幕附近用 speaker classification 与逐帧分类对比，说明有些任务可以减少 query，有些任务不行。}

这点容易被忽略。减少 key/value 主要是减少每个 query 可看的对象；减少 query 则改变要输出多少个结果。前者更多是内部计算优化，后者可能改变任务定义。

\subsection{Compressed Attention 与 Linformer}

课程中提到两种产生代表性 key 的方式。Compressed Attention 可以用卷积等方法把长序列压成短序列；Linformer 则把 key 矩阵乘上一个 $N\times K$ 的投影矩阵。若 key 矩阵写成 $K_{\mathrm{full}}\in\mathbb{R}^{d\times N}$，Linformer 型投影可写成
\[
    K_{\mathrm{compressed}} = K_{\mathrm{full}} E,\qquad E\in\mathbb{R}^{N\times K},
\]
于是 $K_{\mathrm{compressed}}\in\mathbb{R}^{d\times K}$。value 也可用类似方式压缩。

\begin{figure}[H]
    \centering
    \includegraphics[width=0.86\textwidth]{figures/fig11_reduce_keys_linformer.jpg}
    \caption{Compressed Attention 可用卷积压缩序列，Linformer 可用 $N\times K$ 投影矩阵把 $d\times N$ 的 key 压成 $d\times K$。\protect\footnotemark}
    \label{fig:linformer}
\end{figure}
\footnotetext{视频画面时间区间：约 36:44--37:40；字幕附近说明可用 CNN 把长序列变短，或把 $d\times N$ key 矩阵乘上 $N\times K$ 矩阵得到 $d\times K$。}

\begin{importantbox}{Key/Value 压缩的关键问题}
    压缩是否有效，取决于代表性 key/value 是否保留了任务需要的信息。若 $K$ 太小，模型会漏信息；若 $K$ 太大，节省的计算有限。低秩方法通常需要在速度、记忆体和近似误差之间取平衡。
\end{importantbox}

\subsection{本章小结}

低秩和代表性 key/value 方法不是问「哪些格子不算」，而是问「矩阵是否本来就可以更小」。减少 key/value 能降低 attention matrix 的宽度，同时保持输出长度；减少 query 则会改变输出长度，需要看任务是否允许。Linformer 的核心是用线性投影把 $N$ 个 key/value 压成 $K$ 个代表性向量。

% =====================================================================
\section{把 Attention 看成三矩阵乘法}
% =====================================================================

接下来课程进入更代数化的路线。暂时忽略 softmax，self-attention 可以被看成三个矩阵相乘：
\[
    O \approx V K^\top Q.
\]
这里采用课程投影片中的方向，$Q,K\in\mathbb{R}^{d\times N}$，$V\in\mathbb{R}^{d'\times N}$，因此 $K^\top Q\in\mathbb{R}^{N\times N}$，再乘上 $V$ 得到 $O\in\mathbb{R}^{d'\times N}$。

\begin{figure}[H]
    \centering
    \includegraphics[width=0.88\textwidth]{figures/fig12_three_matrix_multiplication.jpg}
    \caption{忽略 softmax 时，attention 可被看成 $V,K^\top,Q$ 三个矩阵相乘；标准做法会先形成 $K^\top Q$ 这个 $N\times N$ 中间矩阵。\protect\footnotemark}
    \label{fig:three-matrix}
\end{figure}
\footnotetext{视频画面时间区间：约 40:45--41:47；字幕附近说明真正做的是 $V$、$K^\top$、$Q$ 三个矩阵相乘，先把 $K^\top$ 和 $Q$ 乘起来，再乘 $V$ 得到 $O$。}

\subsection{乘法顺序为什么重要}

矩阵乘法满足结合律：
\[
    V(K^\top Q)=(VK^\top)Q.
\]
数学上结果相同，但计算成本不同。若先算 $K^\top Q$，中间会产生 $N\times N$ 矩阵；若先算 $VK^\top$，中间矩阵大小是 $d'\times d$。当 $N$ 很大、$d,d'$ 相对较小时，后者显然便宜。

粗略比较如下。若先算 $K^\top Q$：
\[
    K^\top Q: O(N^2d),\qquad
    V(K^\top Q): O(d'N^2).
\]
若改成先算 $VK^\top$：
\[
    VK^\top: O(d'dN),\qquad
    (VK^\top)Q: O(d'dN).
\]
所以后者大约随 $N$ 线性成长，而不是随 $N^2$ 成长。

\begin{warningbox}{为什么不能直接宣布问题解决}
    标准 attention 中间还有 softmax。Softmax 是对每个 query 的所有 key 分数做归一化，它依赖完整的一行 $K^\top Q$。因此只靠矩阵结合律，只能解释「没有 softmax 时」为什么可以加速。若要处理真实 attention，还必须说明 softmax 或类似归一化如何被改写或近似。
\end{warningbox}

\subsection{本章小结}

把 attention 写成三矩阵乘法后，瓶颈变得更清楚：标准顺序会显式生成 $N\times N$ 的 $K^\top Q$；若能合法地改为先算 $VK^\top$，复杂度就可能从平方级降到线性级。接下来的 linear attention 就是在尝试把 softmax 也纳入这种重排。

% =====================================================================
\section{Linear Attention：把 Softmax 拆成特征映射}
% =====================================================================

课程接下来把 softmax 放回公式中。对第一个位置的输出 $b^1$，标准 attention 可以写成
\[
    b^1
    =
    \sum_{i=1}^{N}
    \frac{\exp(q^1\cdot k^i)}
         {\sum_{j=1}^{N}\exp(q^1\cdot k^j)}
    v^i.
\]
这个式子显示，$b^1$ 是所有 value 的 weighted sum，权重来自 $q^1$ 与每个 $k^i$ 的 softmax。

\begin{figure}[H]
    \centering
    \includegraphics[width=0.9\textwidth]{figures/fig13_softmax_attention_formula.jpg}
    \caption{标准 attention 中，$b^1$ 由 $q^1$ 对所有 key 的 softmax 权重加权所有 value 得到。\protect\footnotemark}
    \label{fig:softmax-formula}
\end{figure}
\footnotetext{视频画面时间区间：约 46:44--47:43；字幕附近说明若要计算 $b^1$，就把 $q^1$ 与所有 key 算 attention，再对 value 做 weighted sum。}

\subsection{Kernel 化的关键假设}

Linear attention 的关键一步，是希望把
\[
    \exp(q\cdot k)
\]
改写或近似为两个特征向量的内积：
\[
    \exp(q\cdot k)\approx \phi(q)^\top\phi(k).
\]
这里 $\phi(\cdot)$ 是某种 feature map。不同论文对 $\phi$ 的选择不同：有的使用确定性映射，有的使用随机特征，有的追求无偏估计，有的追求数值稳定。

\begin{figure}[H]
    \centering
    \includegraphics[width=0.9\textwidth]{figures/fig14_kernel_feature_map.jpg}
    \caption{若能把 $\exp(q\cdot k)$ 近似为 $\phi(q)^\top\phi(k)$，就可以把 softmax attention 改写成 feature map 之间的内积形式。\protect\footnotemark}
    \label{fig:kernel-map}
\end{figure}
\footnotetext{视频画面时间区间：约 50:00--50:57；字幕附近开始整理分子分母，并把 $\exp(q\cdot k)$ 替换成 $\phi(q)$ 与 $\phi(k)$ 的内积。}

将近似代入 $b^1$：
\[
    b^1
    =
    \frac{\sum_{i=1}^{N}\phi(q^1)^\top\phi(k^i)v^i}
         {\sum_{j=1}^{N}\phi(q^1)^\top\phi(k^j)}.
\]
因为 $\phi(q^1)$ 与求和索引 $i,j$ 无关，可以把它从求和中提出或重排。分母可写成
\[
    \phi(q^1)^\top\left(\sum_{j=1}^{N}\phi(k^j)\right).
\]
分子则可以看成先把每个 $\phi(k^i)$ 与 $v^i$ 做外积，再对 $i$ 求和，最后由 $\phi(q^1)$ 选择组合。

\subsection{分子重排：从逐位置加权到全局模板}

投影片把分子展开，说明它不是魔法，而是普通代数整理。假设 $\phi(q)$ 和 $\phi(k)$ 的维度为 $M$。那么
\[
    \phi(q^1)^\top\phi(k^i)
    =
    \sum_{m=1}^{M}\phi_m(q^1)\phi_m(k^i).
\]
代回分子：
\[
    \sum_{i=1}^{N}\left(\sum_{m=1}^{M}\phi_m(q^1)\phi_m(k^i)\right)v^i
    =
    \sum_{m=1}^{M}\phi_m(q^1)
    \left(\sum_{i=1}^{N}\phi_m(k^i)v^i\right).
\]
括号中的项
\[
    T_m=\sum_{i=1}^{N}\phi_m(k^i)v^i
\]
只和所有 key/value 有关，不和当前 query 是第几个位置有关。

\begin{figure}[H]
    \centering
    \includegraphics[width=0.94\textwidth]{figures/fig15_numerator_rearrangement.jpg}
    \caption{把分子展开后，可以把乘上同一维 query feature 的项收集在一起；这些括号里的向量只由 key/value 决定，可被多个 query 共用。\protect\footnotemark}
    \label{fig:numerator-rearrangement}
\end{figure}
\footnotetext{视频画面时间区间：约 53:14--54:16；字幕附近说明把乘上 $q^1_1$、$q^1_2$ 的项分别集合，得到若干组 value 的 linear combination。}

从这个角度看，linear attention 先从整个序列中算出一组全局模板 $T_1,\ldots,T_M$，每个位置的 query 再用自己的 $\phi(q^t)$ 去组合这些模板。分母同理，也可以先算出
\[
    z=\sum_{j=1}^{N}\phi(k^j),
\]
再用 $\phi(q^t)^\top z$ 做归一化。

\subsection{为什么它是线性的}

计算 $T_m$ 与 $z$ 都只需要扫过序列一次。对所有位置来说，key/value 产生的这些模板可以共用。计算 $b^1$ 后，计算 $b^2$ 不需要重新求一次所有模板，只需要换成 $\phi(q^2)$ 去组合。

\begin{figure}[H]
    \centering
    \includegraphics[width=0.9\textwidth]{figures/fig16_reuse_templates_linear_attention.jpg}
    \caption{Linear attention 中，由 key/value 产生的蓝色向量和黄色归一化项只需计算一次；不同位置只要换不同的 $\phi(q)$ 去组合它们。\protect\footnotemark}
    \label{fig:reuse-templates}
\end{figure}
\footnotetext{视频画面时间区间：约 57:57--58:39；字幕附近说明计算 $b^1$ 时得到的蓝色 factor 和黄色 factor 在计算 $b^2$ 时不需要重算，只需换成 $\phi(q^2)$。}

若 feature map 维度 $M$ 视为常数或远小于 $N$，整个过程复杂度大约是
\[
    O(NMd')+O(NM),
\]
随序列长度 $N$ 线性成长。代价是：我们不再显式得到完整 attention matrix，也较难像标准 attention 那样直接解释「某个 token attend 到哪个 token」。

\subsection{不同论文如何实现}

课程列出几种代表做法：Efficient Attention、Linear Transformer、Random Feature Attention、Performer。它们都围绕类似思想，但对 $\phi$ 的构造、数值稳定、归一化和随机近似有不同处理。

\begin{figure}[H]
    \centering
    \includegraphics[width=0.88\textwidth]{figures/fig17_linear_attention_realizations.jpg}
    \caption{Efficient Attention、Linear Transformer、Random Feature Attention、Performer 都可被放在线性化 attention 的实现族中；差别主要在 feature map 与近似方式。\protect\footnotemark}
    \label{fig:realizations}
\end{figure}
\footnotetext{视频画面时间区间：约 64:15--64:30；字幕附近说明不同论文用不同方式拆解 $\exp(q\cdot k)$，并列出 Efficient Attention、Linear Transformer、Random Feature Attention、Performer。}

\begin{warningbox}{Linear attention 的理解边界}
    Linear attention 不是把标准 softmax attention 精确免费化。它通常依赖 kernel 近似、feature map 选择或归一化改写。速度变快后，可能牺牲一部分表达能力、训练稳定性或可解释性。实际是否更好，要看任务、序列长度、硬件实现和数值细节。
\end{warningbox}

\subsection{本章小结}

Linear attention 的核心是把 $\exp(q\cdot k)$ 改写为 $\phi(q)^\top\phi(k)$，再利用代数重排，让 key/value 产生的全局模板被所有 query 共用。它绕开显式 $N\times N$ attention matrix，使复杂度有机会随 $N$ 线性成长；但这个收益来自近似和公式改变，而不是无条件等价于标准 attention。

% =====================================================================
\section{还需要用 Q 和 K 算 Attention 吗}
% =====================================================================

前面的方法都仍然保留一个共同设定：attention weight 来自 query 与 key 的关系。最后课程提出更激进的问题：我们真的需要 $Q$ 和 $K$ 吗？

\subsection{Synthesizer：Attention Matrix 可以是参数}

Synthesizer 的想法很反常识。它不一定用 $q$ 和 $k$ 的内积产生 attention matrix。某些版本会让输入经过网络产生矩阵；还有更极端的版本，直接把 attention matrix 视为网络参数的一部分。也就是说，矩阵中的 $N\times N$ 个数本身就是可训练参数。

\begin{figure}[H]
    \centering
    \includegraphics[width=0.86\textwidth]{figures/fig18_synthesizer_parameters.jpg}
    \caption{Synthesizer 可以不由 $Q$ 和 $K$ 计算 attention matrix，而是把矩阵看成网络参数，直接学习一组 mixing weights。\protect\footnotemark}
    \label{fig:synthesizer}
\end{figure}
\footnotetext{视频画面时间区间：约 65:44--66:44；字幕附近说明 Synthesizer 中 attention matrix 不是由 $Q$ 和 $K$ 产生，甚至可以直接作为网络中的 $N\times N$ 个参数。}

这个方法的启示是：self-attention 的本质也许不只是「相似度搜索」，还可以被理解为一种 token mixing。若任务只需要在位置之间混合信息，那么混合矩阵未必一定要来自 query-key dot product。

当然，这也带来明显限制。若 attention matrix 是固定参数，它对输入内容的适应性较弱；序列长度变化时也需要额外处理。它更像是在提醒我们：Transformer 的成功不一定只来自某一个固定解释，许多机制都可能在不同层面贡献效果。

\subsection{Attention-free 与 MLP Mixing}

课程也提到 FNet、Pay Attention to MLPs、MLP-Mixer 等 attention-free 方向。这些方法尝试不用传统 attention，而用 Fourier transform、MLP 或其他 mixing 机制在 token 间传递信息。

它们和 Synthesizer 有相似精神：如果目标是让不同位置交换信息，那么并非只有 $QK^\top$ 一条路。不同 mixing 机制在速度、归纳偏置、硬件友好性和表达能力上各有取舍。

\subsection{本章小结}

Synthesizer 与 attention-free 方法把问题推到更高层：我们到底需要的是 attention 这个公式，还是需要 token 间的信息混合能力？它们不一定取代标准 attention，但能帮助我们理解 Transformer 架构中哪些部分是核心、哪些部分可以被替换。

% =====================================================================
\section{总结与延伸}
% =====================================================================

本讲最后用 Long Range Arena 的结果比较不同方法。图中纵轴代表效果分数，横轴代表速度，圆圈大小代表记忆体使用。可以看到，方法之间没有单一赢家：有些方法很快但效果掉得多，有些方法效果接近但速度优势有限，有些方法记忆体较省但实现复杂。

\begin{figure}[H]
    \centering
    \includegraphics[width=0.9\textwidth]{figures/fig19_lra_summary.jpg}
    \caption{Long Range Arena 比较了多种高效 self-attention 方法：横轴越右速度越快，纵轴越高效果越好，圆圈越大表示记忆体使用越多。\protect\footnotemark}
    \label{fig:lra-summary}
\end{figure}
\footnotetext{视频画面时间区间：约 68:01--70:16；字幕附近说明 LRA score 越高代表表现越好，横轴代表速度，圆圈大小代表 memory，并比较 Transformer、Local Attention、Sinkhorn、Performer 等方法。}

\subsection{一张方法地图}

可以把本讲方法整理为下表。

\begin{table}[H]
    \centering
    \begin{tabular}{@{}L{0.2\textwidth}L{0.32\textwidth}L{0.34\textwidth}@{}}
        \toprule
        路线 & 代表方法 & 核心取舍 \\
        \midrule
        人类设计稀疏图 & Local Attention、Longformer、BigBird & 简单高效，但连接 pattern 固定，可能漏掉输入相关的重要依赖。 \\
        数据相关稀疏化 & Reformer、Routing Transformer & 根据 query/key 相似性筛选连接，适应性更强，但筛选过程要足够便宜。 \\
        可学习 pattern & Sinkhorn Sorting Network & 让网络学习 mask，表达更灵活，但训练和可微近似更复杂。 \\
        低秩压缩 & Linformer、Compressed Attention & 减少 key/value 数量，保持输出长度，但依赖低秩或代表性假设。 \\
        线性化 attention & Linear Transformer、Performer、Efficient Attention & 用 feature map 和代数重排绕开 $N^2$，但通常涉及近似与数值稳定问题。 \\
        新 mixing 框架 & Synthesizer、FNet、MLP-Mixer & 不再执着于 $QK^\top$，强调 token mixing，但归纳偏置和泛化表现要重新评估。 \\
        \bottomrule
    \end{tabular}
    \caption{本讲 self-attention 变型的主要路线。}
    \label{tab:method-map}
\end{table}

\subsection{理解这些方法时最重要的三个问题}

第一，\textbf{它省掉了什么？} 有的方法省掉 attention matrix 中的格子，有的方法省掉 key/value，有的方法省掉显式矩阵，有的方法省掉 query-key 相似度本身。不同「省法」会带来不同副作用。

第二，\textbf{它依赖什么假设？} Local attention 依赖局部性，BigBird 依赖稀疏图仍能保留传播路径，Reformer 依赖相似向量应在同一 cluster，Linformer 依赖低秩，linear attention 依赖 kernel 近似或 feature map 设计。若假设和任务不匹配，速度提升可能换来明显性能损失。

第三，\textbf{它在真实系统中是否真的快？} 理论复杂度只是第一步。GPU/TPU 上的矩阵乘法、稀疏 kernel、memory bandwidth、batch size、序列长度、实现成熟度都会影响最终速度。某些方法理论上更低阶，但如果实现无法充分利用硬件，实际速度未必占优。

\begin{importantbox}{课程给出的总判断}
    Efficient attention 不是一个单一技巧，而是一组围绕长序列瓶颈展开的设计空间。原始 Transformer 很强，但 $N^2$ 成本迫使我们思考：哪些连接必须精算，哪些信息可以压缩，哪些公式可以重写，哪些架构假设可以被替换。
\end{importantbox}

\subsection{延伸阅读方向}

若继续深入，可以按三条线阅读。第一条线是长上下文 Transformer，重点看 Longformer、BigBird、Reformer、Performer 等模型如何在真实长序列任务中取舍。第二条线是 kernelized attention，关注 $\phi(q)^\top\phi(k)$ 如何近似 softmax kernel，以及随机特征如何控制误差和稳定性。第三条线是 token mixing 架构，比较 attention、MLP、Fourier transform、卷积和状态空间模型在序列建模中的角色。

本讲最值得带走的不是某个方法名称，而是一种审题方式：当一个模型模块太贵时，先找出真正贵在哪里，再问它是否必须完整计算。Self-attention 的瓶颈来自 $N\times N$ 两两交互，而本讲所有方法都在用不同方式回答「哪些交互可以省、压、近似或替换」。

\end{document}
